\begin{figure}[htbp]
\centering
\begin{tikzpicture}[x=1.55cm,y=1.55cm,font=\small,>=Latex]
 \pgfmathsetmacro{\bandheight}{sqrt(3)/2}
 \fill[blue!4] (-.8,\bandheight) rectangle (3.2,2*\bandheight);
 \begin{scope}
  \clip (-.8,0) rectangle (3.2,3*\bandheight);
  \foreach \j in {0,1,2,3}{
   \draw[black!22] (-1,\j*\bandheight)--(4,\j*\bandheight);
  }
  \foreach \j in {0,1,2}{
   \foreach \i in {-3,-2,...,4}{
    \draw[black!22] ({\i+\j/2},\j*\bandheight)
      --({\i+\j/2+.5},{(\j+1)*\bandheight});
    \draw[black!22] ({\i+\j/2},\j*\bandheight)
      --({\i+\j/2-.5},{(\j+1)*\bandheight});
   }
  }
 \end{scope}
 \draw[thick] (-.8,0)--(3.2,0);
 \draw[thick] (-.8,3*\bandheight)--(3.2,3*\bandheight);
 \coordinate (a) at (.5,0);
 \coordinate (v1) at (.5,\bandheight/3);
 \coordinate (v2) at (1,2*\bandheight/3);
 \coordinate (v3) at (1,4*\bandheight/3);
 \coordinate (v4) at (.5,5*\bandheight/3);
 \coordinate (v5) at (.5,7*\bandheight/3);
 \coordinate (v6) at (1,8*\bandheight/3);
 \coordinate (b) at (1,3*\bandheight);
 \draw[very thick,blue!65!black] (a)--(v1)--(v2)--(v3)--(v4)--(v5)--(v6)--(b);
 \foreach \v in {v1,v2,v3,v4,v5,v6}{\fill (\v) circle (1.4pt);}
 \foreach \v in {a,b}{\draw[thick,fill=white] (\v) circle (2pt);}
 \node[below=3pt] at (a) {initial port};
 \node[above=3pt] at (b) {terminal port};
 \node[align=left,anchor=west] (vertex) at (1.65,1.05) {visited dual\\vertex};
 \draw[->] (vertex.west)--(v3);
 \draw[<->] (3.4,\bandheight)--node[right]{one band}(3.4,2*\bandheight);
 \draw[black!50] (3.2,\bandheight)--(3.5,\bandheight);
 \draw[black!50] (3.2,2*\bandheight)--(3.5,2*\bandheight);
\end{tikzpicture}
\caption{Port and length conventions. The path passes through six adjacent
triangles of the triangular tiling, so it visits six vertices of the
honeycomb dual and has port weight $\rho^6$. Its two terminal pieces are
half-edges. An ordinary vertex-to-vertex path counts full edges instead. The shaded region is one triangular band.}
\label{fig:model-conventions}
\end{figure}
